%% ============================================================================
\section{Background}
\label{sec:background}
%% ============================================================================

\subsection{Lubrication regimes and the viscosity ratio $\kappa$}
\label{subsec:kappa}

Ball bearings operate in one of three lubrication regimes depending on the extent to which the lubricant film separates the rolling surfaces. In the full-film elastohydrodynamic lubrication (EHL) regime the surfaces are completely separated and asperity contact is suppressed; in the mixed regime the film is partially formed and intermittent asperity contact occurs; in the boundary regime the film is insufficient and metal-to-metal contact dominates. The ISO 281 viscosity ratio $\kappa$ provides a continuous, dimensionless indicator of the lubricant viscosity relative to that required for full EHL film formation. It is commonly used for lubrication regime classification with approximate boundaries at $\kappa < 0.1$ (boundary), $0.1 \leq \kappa < 0.4$ (mixed, severe),  $0.4 \leq \kappa < 1$ (mixed, moderate), and $\kappa \geq 1$ (full film) \cite{schaeffler_lubrication_2013,skf_viscosity_ratio}. The ISO~281 viscosity ratio is defined as
\begin{equation}
  \kappa(T, n, d_{pw}) = \frac{\nuop(T)}{\nureq(n, d_{pw})},
  \label{eq:kappa}
\end{equation}
where $\nuop(T)$ [\si{mm^2\,s^{-1}}] is the actual kinematic viscosity of the lubricant at operating temperature $T$, and $\nureq(n, d_{pw})$ is the minimum kinematic viscosity required for full EHL at the bearing's rotational speed $n$ [\si{rpm}] and pitch diameter $d_{pw}$ [\si{mm}]~\cite{iso_281}. The reference viscosity is given by
\begin{equation}
  \nureq(n, d_{pw}) = \begin{cases}
    45000n^{-0.83}d_{pw}^{-0.5} & n < \SI{1000}{rpm}, \\
    4500n^{-0.5}d_{pw}^{-0.5} & n \geq \SI{1000}{rpm}.
  \end{cases}
  \label{eq:nu1}
\end{equation}
The operating viscosity $\nuop$ is computed from the ASTM~D341 Walther viscosity--temperature equation~\cite{astm_d341} fitted to the lubricant's \SI{40}{\celsius} and \SI{100}{\celsius} reference values, with a standard value of $\lambda=0.7$. Since $\nureq$ decreases with increasing rotational speed, $\kappa$ captures the combined effect of speed, temperature, and lubricant grade on lubrication adequacy at a given operating point. Importantly, $\kappa$ assumes nominally intact lubricant and sufficient lubricant supply: the Walther equation reflects the properties of fresh lubricant, so $\kappa$ does not encode viscosity changes from degradation or contamination, nor does it account for starvation. At high rotational speeds, replenishment may not keep pace with lubricant displaced by the rolling elements, reducing the actual film below $\kappa$'s prediction even when the bulk lubricant is intact~\cite{poll_starved_2019}.

\subsection{Acoustic emission and passive ultrasound for lubrication monitoring}
\label{subsec:origins}
AE refers to transient elastic stress waves generated by sudden energy releases in a material, where asperity contact and micro-fracture are the dominant sources in rolling bearings~\cite{renhart_ae_2021,miettinen_analysis_2001}. Under boundary and mixed lubrication, intermittent asperity contacts generate broadband stress waves that propagate through the bearing structure and are detected by AE sensors coupled to the outer housing~\cite{cornel_condition_2021,renhart_ae_2021}. The emission rate, amplitude, and frequency content of these events are therefore sensitive to changes in contact severity, with intensity decreasing as film adequacy improves. Empirically, AE energy and pulse-count measures track lubrication condition \cite{miettinen_analysis_2001,yoshioka_compound_2009,hou_insitu_2025}, AE shows earlier sensitivity to contamination and faults than vibration \cite{tandon_ae_2006,cornel_condition_2021,wang_fingerprint_2023}, and its frequency content and power have been tied to film thickness through data-driven \cite{renhart_ae_2021} and physical models \cite{liu_combined_2021,zhang_coupled_2026}. Transmission path loss is inhomogeneous and path-dependent, affecting both the amplitude and frequency content of the received signal~\cite{vallen_sensor_preamplifier_2021}. AE sensors for bearings respond from tens of kHz to several hundred kHz.

Passive US sensing detects structure-borne acoustic energy emanating from the bearing during operation and has been explored far less than AE for lubrication monitoring outside commercial threshold tools \cite{jakobsen_detecting_2023}. Ultrasonic emissions span from 20~kHz into the MHz range, overlapping much of the AE band. The level and spectral character of US signals are sensitive to the lubrication condition, friction level, and asperity contact activity~\cite{jakobsen_detecting_2023,lineham_ultrasonic_2008}. Commercial handheld probes (UE Systems, Sonotec, SDT Ultrasound) exploit this by comparing the signal level against a baseline and triggering relubrication when a threshold is exceeded~\cite{lineham_ultrasonic_2008}, making threshold-based passive US a practical non-invasive LCM approach in routine industrial maintenance. These probes typically use heterodyned US down-mixed to the audible range to enable human-assisted diagnostics; the sensors are narrowband, with peak sensitivity at their resonance frequency.

The two sensing methods thus respond to the same underlying contact events---asperity contacts, rapid pressure changes, friction, particle impacts, and crack growth---but through different physical mechanisms and at different frequency scales. AE captures individual high-frequency transient events from discrete micro-fractures and asperity collisions, while passive US reflects structural resonances, friction noise, and the aggregate effect of contact activity. Both channels should therefore carry information about $\kappa$ but their different sensitivities to transient events and continuous contact phenomena suggest that AE and passive US may provide complementary information about lubrication-dependent contact dynamics. These considerations motivate the experimental investigation that follows, which seeks to determine which of these information sources explains the predictive performance of acoustic $\kappa$-regression models.